\documentclass[11pt,a4paper]{article}
\usepackage[T1]{fontenc}
\usepackage{lmodern,amsmath,amssymb,amsthm,mathtools}
\usepackage[margin=25mm]{geometry}
\usepackage{xcolor}
\usepackage[colorlinks=true,linkcolor=blue!45!black,citecolor=blue!45!black,urlcolor=blue!45!black]{hyperref}
\usepackage{microtype}
\newtheorem{theorem}{Theorem}[section]
\newtheorem{proposition}[theorem]{Proposition}
\newtheorem{lemma}[theorem]{Lemma}
\newtheorem{corollary}[theorem]{Corollary}
\theoremstyle{definition}
\newtheorem{definition}[theorem]{Definition}
\newtheorem{remark}[theorem]{Remark}
\newcommand{\A}{\mathbb A}
\setlength{\emergencystretch}{2em}
\hypersetup{pdftitle={Polynomial systems and the field of witnesses},pdfsubject={Variable-preserving conjunctions and bounded-degree extension obstructions}}
\title{Polynomial Systems and the Field of Witnesses}
\author{}
\date{8 November 2026}
\begin{document}
\maketitle
\begin{abstract}
We combine $s$ polynomial equations over a number field into one
without introducing any coordinates, with degree at most
$\lceil\sqrt{s}\rceil$ times the largest original degree. The
construction uses a homogeneous anisotropic form obtained from
multiplication determinants at one finite prime. Over a formally
real field the familiar sum of squares gives degree at most twice
the largest original degree, independently of $s$. We also prove
that an existential formula with infinitely many inputs in a
number field acquires all but finitely many base-field inputs
after extensions of uniformly bounded degree. For a single
nonzero polynomial this degree bound can be taken to be its smallest
positive individual degree in a witness coordinate. Consequently
one unchanged existential formula cannot define the rings of
integers, or the corresponding rings of $S$-integers, in all finite
extensions of a number field.
\end{abstract}

\section{Combining equations without new coordinates}

For a field $K$, let $F_1,\ldots,F_s\in K[X_1,\ldots,X_n]$,
where $s\ge1$. The common zero set is evaluated on $K^n$.
We seek a polynomial $P$ in the same coordinates satisfying
\[
 P(\mathbf x)=0
 \quad\Longleftrightarrow\quad
 F_1(\mathbf x)=\cdots=F_s(\mathbf x)=0
                    \qquad(\mathbf x\in K^n).
\]
The equivalence concerns points over $K$; an extension of the
field changes the assertion unless its proof also applies there.

\begin{proposition}\label{sys:sum-squares}
If $K$ is formally real, the zero set of $\sum_{i=1}^s F_i^2$
in $K^n$ is the common zero set of $F_1,\ldots,F_s$. Its degree is at most
$2\max_i\deg F_i$.
\end{proposition}
\begin{proof}
Formal reality means that $-1$ is not a sum of squares in $K$.
If a vanishing sum of squares had a nonzero summand, division
by that square would express $-1$ as a sum of the other squares,
a contradiction. Squaring gives the degree bound.
\end{proof}

Every number field with a real embedding is formally real. Over
other number fields an anisotropic form combines the equations.
A homogeneous form $A\in K[Z_1,\ldots,Z_r]$ is
called \emph{anisotropic over $K$} if its only zero in $K^r$ is
the zero vector.

\begin{lemma}\label{sys:anisotropic-form}
Let $K$ be a number field and $e\ge1$ an integer. There is an
explicit homogeneous anisotropic form
\[
       A_e\in\mathcal O_K[Z_{j,i}:0\le j,i<e]
\]
of degree $e$ in $e^2$ coordinates.
\end{lemma}
\begin{proof}
Choose a nonzero prime ideal $\mathfrak p$ of $\mathcal O_K$,
normalize $v_{\mathfrak p}$ to have value group $\mathbb Z$,
and choose an element $\pi\in\mathfrak p\setminus\mathfrak p^2$ and a
monic polynomial
\[
              g(U)\in\mathcal O_K[U],\qquad \deg g=e,
\]
whose reduction modulo $\mathfrak p$ is irreducible. Such a
polynomial is obtained by lifting the coefficients of an
irreducible degree-$e$ polynomial over the finite residue field.

In the $e$-dimensional $K$-algebra $B=K[U]/(g)$, use the basis
$1,u,\ldots,u^{e-1}$, where $u$ is the image of $U$. Define
\[
 N_g(Z_0,\ldots,Z_{e-1})
   =\det\left(m_{\sum_{i=0}^{e-1}Z_i u^i}\right),
\]
where $m_b$ denotes multiplication by $b$ in this basis.
Reduction by the monic relation $g(u)=0$ gives a matrix whose
entries are linear forms in the $Z_i$ with coefficients in
$\mathcal O_K$. Thus $N_g$ is a homogeneous integral form of
degree $e$, with $N_g(1,0,\ldots,0)=1$.

Write $v=v_{\mathfrak p}$. For any nonzero vector
$\mathbf z\in K^e$, put $a=\min_i v(z_i)$. The entries of
$\pi^{-a}\mathbf z$ belong to the valuation ring and are not
all in its maximal ideal. Modulo that ideal, the multiplication
matrix is multiplication by a nonzero element of the field
\[
          (\mathcal O_K/\mathfrak p)[U]/(\overline g).
\]
Its determinant is nonzero in the residue field. Homogeneity
therefore proves the exact valuation identity
\begin{equation}\label{sys:norm-valuation}
                    v(N_g(\mathbf z))=e\min_i v(z_i).
\end{equation}
In particular, $N_g(\mathbf z)\ne0$ for every nonzero
$\mathbf z\in K^e$.

Now set
\begin{equation}\label{sys:weighted-norm-form}
 A_e((Z_{j,i}))
    =\sum_{j=0}^{e-1}\pi^j N_g(Z_{j,0},\ldots,Z_{j,e-1}).
\end{equation}
At a nonzero input, omit any zero blocks. By
\eqref{sys:norm-valuation}, the valuations of the remaining
summands are congruent to their distinct indices $j$ modulo $e$.
Exactly one summand has least valuation. The sum cannot vanish.
This proves anisotropy. Homogeneity and the coefficient assertion
follow directly from \eqref{sys:weighted-norm-form}.
\end{proof}

\begin{theorem}\label{sys:combine-field}
Let $K$ be a number field and let $F_1,\ldots,F_s$ be polynomials
in a fixed set of coordinates, each of degree at most $d$.
Put $e=\lceil\sqrt{s}\rceil$. Padding the list of equations with
zeros to length $e^2$, the polynomial
\[
                     P=A_e(F_1,\ldots,F_s,0,\ldots,0)
\]
has the same $K$-rational zero set as the original system, uses
the same coordinates, and has degree at most $ed$.
If the $F_i$ have coefficients in
$\mathcal O_K$, then so does $P$.
\end{theorem}
\begin{proof}
The form $A_e$ vanishes only at the zero vector over $K$.
Substituting the $F_i$ therefore gives the stated equivalence.
Since $A_e$ is homogeneous of degree $e$, each resulting term
has degree at most $ed$. No coordinate is introduced by the
substitution.
\end{proof}

For example, any four equations of degree at most $d$ over a
number field admit one equivalent equation of degree at most $2d$, and
any nine admit one of degree at most $3d$. For a formally real
number field, Proposition~\ref{sys:sum-squares} gives the
stronger bound $2d$ for arbitrary $s$. For $s=1$ one uses the
original polynomial.

\subsection*{A concrete quadratic form}

Over $K=\mathbb Q(i)$, take the prime ideal $(3)$, whose residue
field has nine elements. The element $1+i$ is not a square in
that field, since $(1+i)^4=-1\ne1$ modulo $3$. Applying the
construction with $e=2$, $\pi=3$ and $g(U)=U^2-(1+i)$ gives
\[
 A_2(Z_1,Z_2,Z_3,Z_4)
     =Z_1^2-(1+i)Z_2^2+3Z_3^2-3(1+i)Z_4^2.
\]
It is anisotropic over $\mathbb Q(i)$ by the valuation proof
above. It combines four equations into one even though
a plain sum of squares is unsuitable over that field.

For a cubic form in nine coordinates over the same field, use
$\mathfrak p=(1+i)$, $\pi=1+i$ and $g(U)=U^3+U+1$.
The reduction of $g$ has no root in $\mathbb F_2$ and is
therefore irreducible. The multiplication determinant is
\begin{align*}
 N(a,b,c)
  &=\det\begin{pmatrix}
       a&-c&-b\\ b&a-c&-b-c\\ c&b&a-c
       \end{pmatrix}\\
  &=a^3-2a^2c+ab^2+3abc+ac^2-b^3-bc^2+c^3.
\end{align*}
Consequently
\[
 N(Z_1,Z_2,Z_3)+(1+i)N(Z_4,Z_5,Z_6)
                      +2iN(Z_7,Z_8,Z_9)
\]
is an explicit anisotropic cubic over $\mathbb Q(i)$. Its
coefficient list and the determinant identity are recorded in
the accompanying exact computational data.

\subsection*{Comparison with one norm}

One can also choose an extension $L/K$ of degree $s$ and a
$K$-basis $b_1,\ldots,b_s$. Then
\[
              N_{L/K}\left(\sum_{i=1}^s b_iF_i\right)=0
\]
is equivalent to all $F_i=0$, because the norm of a nonzero
field element is nonzero. This gives degree at most $sd$.
Theorem~\ref{sys:combine-field} obtains the square-root bound by
combining several smaller multiplication determinants at
distinct valuation classes.

\section{Existential formulas after finite extensions}

Fix a number field $k$. An existential formula considered here
has the form
\[
 \Phi(T):\quad \exists Y_1,\ldots,Y_m\quad
             \theta(T,Y_1,\ldots,Y_m),
\]
where $\theta$ is a finite Boolean combination of polynomial
equations with coefficients in $k$. Equivalently, it is a finite
disjunction of conjunctions of equations and nonzero conditions.
For each extension field $L/k$, write
\[
                \Phi(L)=\{a\in L:\Phi(a)\text{ holds in }L\}.
\]
The same formula and the same coefficients are used for every $L$.

\begin{theorem}\label{sys:bounded-extension}
If $\Phi(k)$ is infinite, then there exist a positive integer
$D$ and a nonzero polynomial $h\in k[T]$ such that every
$a\in k$ with $h(a)\ne0$ belongs to $\Phi(L_a)$ for some finite
extension $L_a/k$ of degree at most $D$.
\end{theorem}
\begin{proof}
The conditions defining the witness tuples give a constructible
subset of affine space. By Chevalley's theorem its image in
the input line is constructible \cite{StacksChevalley}. Since
the image has infinitely many $k$-points, it contains the
generic point of the line. Equivalently, at least one of the
finitely many conjunctions has a nonempty generic fiber over
$k(T)$.

Choose a closed point in a nonempty affine open of that generic
fiber. Its residue field $E$ is a finite extension of $k(T)$.
The coordinates of the point give elements
$y_1,\ldots,y_m\in E$ satisfying one conjunction, with every
required nonzero expression nonzero in $E$. Characteristic zero
gives a primitive element $\alpha$ for $E/k(T)$. Let
$G(T,Z)\in k(T)[Z]$ be its monic minimal polynomial, of
degree $D=[E:k(T)]$ in $Z$.

Express each $y_i$, and also the inverse of each required
nonzero expression, as a polynomial in $\alpha$ of degree
less than $D$ with coefficients in $k(T)$. Choose a nonzero
$h\in k[T]$ which clears all denominators in these expressions
and in $G$. Put $R=k[T,1/h]$ and
\[
                              C=R[Z]/(G).
\]
Because $G$ is monic, $C$ is free of rank $D$ over $R$ and
embeds in $E$ under $Z\mapsto\alpha$. The chosen coordinate
expressions and inverse expressions belong to $C$. All the
equations, and the identities asserting that the chosen inverses
are inverses, hold in $C$: they hold in $E$, and $C\hookrightarrow E$
is injective.

Now let $a\in k$ with $h(a)\ne0$. Specialization gives the
nonzero $D$-dimensional algebra
\[
                             C_a=k[Z]/(G(a,Z)).
\]
Choose a maximal ideal of $C_a$. Its residue field $L_a$ is a
finite extension of $k$ of degree at most $D$. The coordinate
expressions descend to a witness tuple in $L_a$; all equations
and inverse identities survive. Hence the same conjunction,
and therefore $\Phi(a)$, holds in $L_a$.
\end{proof}

The finite exceptional set in the theorem consists of the
$k$-roots of the single fixed polynomial $h$. The degree bound
$D$ is uniform in $a$, although the extension $L_a$ may depend
on $a$. The proof specifies $D$ as the degree of one generic
closed point; it does not identify one extension that contains
witnesses for all inputs.

\begin{corollary}\label{sys:no-uniform-integers}
No single existential formula over a number field $k$ defines
$\mathcal O_L$ in every finite extension $L/k$. More precisely,
if $\Phi(k)$ is infinite, there is a positive integer $D$ such
that, for all but finitely many $q\in\mathbb Q\setminus\mathbb Z$,
the formula accepts $q$ in some extension $L_q/k$ of degree at
most $D$.
\end{corollary}
\begin{proof}
Apply Theorem~\ref{sys:bounded-extension} to the elements
$q\in\mathbb Q\subset k$. Every rational algebraic integer is
an integer. Indeed, if $q=a/b$ is in lowest terms and satisfies
a monic equation of degree $r$ with integer coefficients,
clearing denominators yields $b\mid a^r$, so $b=1$.
Consequently each admitted rational noninteger lies outside
$\mathcal O_{L_q}$.
\end{proof}

\begin{corollary}\label{sys:no-uniform-s-integers}
Let $S$ be a fixed finite set of finite primes of $k$, and for
each finite extension $L/k$ let $S_L$ be the primes of $L$
above $S$. One unchanged existential formula over $k$ cannot
define $\mathcal O_{L,S_L}$ in every finite extension $L/k$.
\end{corollary}
\begin{proof}
Such a formula would have infinitely many inputs over $k$,
since $\mathbb Z\subseteq\mathcal O_{k,S}$. Choose a rational
prime $p$ lying below no prime in $S$ and such that $h(1/p)\ne0$
for the polynomial in Theorem~\ref{sys:bounded-extension}.
There are infinitely many eligible primes and only finitely
many roots of $h$. The theorem admits $1/p$ over some finite
extension $L/k$. At every prime of $L$ above $p$ this element
has negative valuation, and none of those primes belongs to
$S_L$. Thus $1/p\notin\mathcal O_{L,S_L}$.
\end{proof}

These corollaries concern a fixed formula whose witness field
is enlarged. They leave open the separate possibility of a
different defining formula for each number field.

\section{A sharper extension bound for one polynomial}

\begin{proposition}\label{sys:hypersurface-extension-bound}
Let $k$ be an infinite field and
$P\in k[T,Y_1,\ldots,Y_m]$. Suppose that $P$ depends on the
witness coordinate $Y_j$, and put $\delta=\deg_{Y_j}P\ge1$.
There is a nonzero $h\in k[T]$ such that for every $a\in k$
with $h(a)\ne0$, the equation $P(a,\mathbf y)=0$ has a solution
over an extension of $k$ of degree at most $\delta$.
\end{proposition}
\begin{proof}
Regard $P$ as a polynomial in $Y_j$. Its leading coefficient
is a nonzero polynomial in $T$ and the remaining witness
coordinates. Since $k$ is infinite, choose constants for those
remaining coordinates so that the leading coefficient becomes
a nonzero polynomial $h(T)$. The specialization is a polynomial
$Q(T,Z)\in k[T,Z]$ with $\deg_ZQ=\delta$. For $h(a)\ne0$,
choose any irreducible factor of the degree-$\delta$ polynomial
$Q(a,Z)$. Adjoining a root gives an extension of degree at
most $\delta$ and a witness for the original equation.
\end{proof}

\begin{corollary}\label{sys:single-polynomial-extension}
Let $k$ be a number field, and suppose a single nonzero polynomial
$P(T,\mathbf Y)$ has an infinite input projection over $k$.
Then $P$ depends on at least one witness coordinate. In
Corollary~\ref{sys:no-uniform-integers}, the bound on the extension
degree can be taken to be
\[
            D=\min_{j:\,\deg_{Y_j}P>0}\deg_{Y_j}P
               \le \deg P.
\]
In particular an unchanged quartic equation cannot define the
rings of integers over every extension of degree at most four.
If it has an essential witness of individual degree two, it
already fails over an extension of degree at most two.
\end{corollary}
\begin{proof}
If $P$ were independent of all witness coordinates, its input
projection would be the finite zero set of a nonzero polynomial
in $T$. Choose an essential coordinate of least positive
individual degree, apply Proposition~\ref{sys:hypersurface-extension-bound},
and take any rational noninteger outside its finite exceptional
set. The same argument applies to all such rational nonintegers.
\end{proof}

When the polynomial already defines an integral subset over
$k$, the failure in the last sentence cannot occur over $k$
itself; a witnessing extension of degree at most two must then
have degree exactly two. This is an obstruction to transferring
the equation unchanged, rather than a counterexample to
field-dependent Diophantine definability.

\subsection*{An explicit quadratic extension}

Let $P_{392}$ be the quartic specified by the accompanying complete
coefficient and residual tables. Coordinates are numbered $1,\ldots,392$,
with input $X_1$. It is the sum of squares of $364$ integer quadratic
residuals. Put
\[
 A=80989^{54929},\qquad
 D=16P_{392}(\tfrac12,0,\ldots,0).
\]
This formula defines a fixed integer without requiring its decimal
expansion. The exact computation in the supplied coefficient data gives
\[
 D=191779334771151081400576\,A^2+D_0,\qquad D_0>0,
\]
where the complete decimal value of $D_0$ is recorded in
\texttt{field-specialization392.json}.

\begin{proposition}\label{sys:explicit-quartic-extension}
Let $L=\mathbb Q(\theta)$ with $\theta^2=-D$. Set $X_1=1/2$,
$X_{293}=\theta/4$, and every other witness coordinate to zero.
Then $P_{392}=0$ over $L$, but $X_1$ is not an algebraic integer.
Thus this unchanged quartic does not define $\mathcal O_L$ in $L$.
\end{proposition}
\begin{proof}
Leave coordinate $293$ free, calling it $Y$, and substitute the
specified values in the quadratic residuals. Of the $364$ residuals,
$302$ vanish, $61$ become constants, and the residual with source
identifier $493$, the square-output definition, becomes $-Y$.
Every constant belongs to $\tfrac14\mathbb Z[A]$, and one has absolute value $1/4$.
These statements follow by inspection of the explicitly listed
integer terms; the separate checker also verifies them by exact
substitution. Squaring and adding therefore gives
\[
 P_{392}(\tfrac12\mathbf e_1+Y\mathbf e_{293})
                      =Y^2+D/16.
\]
In particular $D$ is a positive integer. The independent checker
obtains this same identity both from all quadratic residuals and from
the complete quartic coefficient table, including its fixed value of $A$.
Since $D>0$, the polynomial $Z^2+D$ is irreducible over $\mathbb Q$,
so $L$ is an imaginary quadratic field. Taking $Y=\theta/4$ makes
the displayed value zero. The rational number $1/2$ is not an
algebraic integer, as proved in Corollary~\ref{sys:no-uniform-integers}.
\end{proof}

\begin{thebibliography}{9}
\bibitem{StacksChevalley}
The Stacks Project Authors, \emph{The Stacks Project},
Chevalley's theorem, Tag 00FE.
\url{https://stacks.math.columbia.edu/tag/00FE}.
\end{thebibliography}
\end{document}
